% APPROVAL ONLY. Do not input this file as a combined manuscript patch.
% A and B are mutually exclusive decisions; only selected fragments may be applied.
% No authoritative current RevThree/RevFive source was located.
% Preserve actual current enclosing revision macros; never strip or invent them.

% section3_weighted_medoid | A: describe traced historical implementation
% Existing source: /home/bardia/projects/graphical-sampling/DI_reference_sensitivity_JSSAM/sources/Measuring_Spatial_Spread_current_2026-10-12.tex:501-509
\[
m_i
\in
\arg\min_{u\in U_i}
\sum_{\ell\in U_i}\pi_\ell
\|\bm c_u-\bm c_\ell\|,
\qquad
\bm r_i=\bm c_{m_i}.
\]

% section3_assignment | A: describe traced historical implementation
% Existing source: /home/bardia/projects/graphical-sampling/DI_reference_sensitivity_JSSAM/sources/Measuring_Spatial_Spread_current_2026-10-12.tex:545-556
The sample units are then matched to the medoid reference points by solving the
assignment problem
\[
\tau
=
\arg\min_{\tau}
\sum_{i=1}^n
\|\bm c_{\tau(i)}-\bm r_i\|,
\qquad
\tau(i)\in s,
\]
where \(\tau\) is a one-to-one mapping from clusters to sampled units.

% section3_historical_kernel | A: describe traced historical implementation
% Existing source: /home/bardia/projects/graphical-sampling/DI_reference_sensitivity_JSSAM/sources/Measuring_Spatial_Spread_current_2026-10-12.tex:586-597
The historical implementation used an isotropic scalar bandwidth
\(h=N^{-1/(d+4)}\), the Scott bandwidth returned by the fitted
\texttt{KernelDensity} estimator, with \(\mathbf H=h^2\mathbf I_d\).
The multivariate uniform kernel is uniform on the Euclidean unit ball:
\[
K(\bm u)=\frac{1}{v_d}\mathbf 1\{\|\bm u\|_2\leq1\},
\qquad
v_d=\frac{\pi^{d/2}}{\Gamma(1+d/2)},
\qquad \bm u\in\mathbb R^d.
\]
Applying the same estimator to the translated coordinates

% section4_replications | A: describe traced historical implementation
% Existing source: /home/bardia/projects/graphical-sampling/DI_reference_sensitivity_JSSAM/sources/Measuring_Spatial_Spread_current_2026-10-12.tex:977-978
\citep{sampling_pkg}. The retained records distinguish Monte Carlo evaluation
from exact enumeration. Where the Monte Carlo replication count can be
recovered, the retained records report \(2{,}000\) samples. Counts remain
unresolved for other blocks, so a common \(5{,}000\)-replication count is not
established for all reported comparisons.

% algorithm2_medoid | A: describe traced historical implementation
% Existing source: /home/bardia/projects/graphical-sampling/DI_reference_sensitivity_JSSAM/sources/Measuring_Spatial_Spread_current_2026-10-12.tex:1490-1505
\State For each cluster \(U_i\), define the reference point \(\bm q_i\) by
\[
\bm q_i
=
\begin{cases}
\bm c_{m_i}, &
\rho=\textsc{medoid},\quad
m_i\in\displaystyle\arg\min_{u\in U_i}
\sum_{\ell\in U_i}\pi_\ell\|\bm c_u-\bm c_\ell\|, \\[2ex]
\displaystyle
\frac{\sum_{\ell\in U_i}\pi_\ell\bm c_\ell}
{\sum_{\ell\in U_i}\pi_\ell}, &
\rho=\textsc{centroid}.
\end{cases}
\]
For the medoid option, ties are broken by largest IP and then smallest label.

% algorithm2_assignment | A: describe traced historical implementation
% Existing source: /home/bardia/projects/graphical-sampling/DI_reference_sensitivity_JSSAM/sources/Measuring_Spatial_Spread_current_2026-10-12.tex:1508-1518
\State Match the reference points to the sampled units by solving
\[
\tau
=
\arg\min_{\tau}
\sum_{i=1}^n
\|\bm c_{\tau(i)}-\bm q_i\|,
\qquad
\tau(i)\in s,
\]
where \(\tau\) is one-to-one.

% comment5_translation_qualification | Both A and B
% Existing source: /home/bardia/projects/graphical-sampling/DI_reference_sensitivity_JSSAM/sources/Measuring_Spatial_Spread_current_2026-10-12.tex:901-917
    \item[(iii)] Conditional on a fixed partition, fixed reference points,
and a fixed matching, DI is invariant to a common translation of the matched
sample coordinates, provided the density estimator is translation equivariant.
Indeed, replacing each matched sample coordinate by
\(\bm c_{\tau(i)}+\bm a\) replaces every translated population coordinate
by \(\widetilde{\bm c}_\ell+\bm a\), leaving its density at the corresponding
translated location unchanged. In particular, if
\(\bm c_{\tau(i)}=\bm r_i+\bm a\) for all \(i\), then the transformed
population is a common translation of the original population and DI is zero.
This conditional property does not imply invariance when the sample alone is
shifted and its sample-adaptive partition, medoids, or matching are recomputed.
With ordinary Euclidean assignment, optimal rematching need not preserve the
original matching. Different spatial spread indices formalize spatial balance
from different viewpoints; DI evaluates departure from the sample-induced
IP\(n\)M reference configuration illustrated in
Figure~\ref{fig:sample_adaptive_reference}. Its rotation properties are discussed
in Appendix~\ref{Exm:rotation_4x5}.

% section3_keep_squared_alternative | B: retain squared definition; explicitly separate historical pipeline
% Existing source: /home/bardia/projects/graphical-sampling/DI_reference_sensitivity_JSSAM/sources/Measuring_Spatial_Spread_current_2026-10-12.tex:545-556
The formal definition below uses squared Euclidean assignment and the
squared-distance representative specified above. The retained historical
results were produced by a distinct pipeline using ordinary Euclidean
assignment, weighted ordinary-distance medoids, and an isotropic spherical
\texttt{tophat} kernel. Those results therefore describe the historical
pipeline and are not numerical validation of the present squared-distance,
axis-aligned-kernel definition.

The sample units are then matched to the medoid reference points by solving the
assignment problem
\[
\tau
=
\arg\min_{\tau}
\sum_{i=1}^n
\|\bm c_{\tau(i)}-\bm r_i\|^2,
\qquad
\tau(i)\in s,
\]
where \(\tau\) is a one-to-one mapping from clusters to sampled units.

% historical_clustering_disclosure | Both A and B; essential historical-method disclosure
% Existing source: /home/bardia/projects/graphical-sampling/DI_reference_sensitivity_JSSAM/sources/Measuring_Spatial_Spread_current_2026-10-12.tex:488-489
\(\{\bm c_\ell:\ell\in s\}\) as initial centers, the formal definition applies
the hard IP\(n\)M procedure to obtain \(n\) spatially compact clusters.
The retained historical results used the then-current weighted \(k\)-means
initialization and refinement path rather than the entropic OT path described
in Section~2; results from that historical implementation are distinguished
from the formal OT definition.

% abstract_translation | Both A and B
% Existing source: /home/bardia/projects/graphical-sampling/DI_reference_sensitivity_JSSAM/sources/Measuring_Spatial_Spread_current_2026-10-12.tex:143-143
bounded measure with a conditional translation property that distinguishes under-dispersion,

% introduction_translation | Both A and B
% Existing source: /home/bardia/projects/graphical-sampling/DI_reference_sensitivity_JSSAM/sources/Measuring_Spatial_Spread_current_2026-10-12.tex:233-235
interpreted as approximately well spread. With the reference partition and matching held fixed, a common translation of the matched sample coordinates leaves DI unchanged. This conditional property concerns the relative arrangement of the matched sample and does not guarantee invariance after sample-adaptive reclustering or rematching.

% conclusion_translation | Both A and B
% Existing source: /home/bardia/projects/graphical-sampling/DI_reference_sensitivity_JSSAM/sources/Measuring_Spatial_Spread_current_2026-10-12.tex:1120-1121
The proposed index is bounded and directly interpretable, and has the conditional translation property described in Section~\ref{Sec:Measure}. Negative values indicate under-dispersion or local clustering,

% appendix_translation_example | Both A and B
% Existing source: /home/bardia/projects/graphical-sampling/DI_reference_sensitivity_JSSAM/sources/Measuring_Spatial_Spread_current_2026-10-12.tex:1150-1151
In panel~(b), a uniform translation is applied to all sample units and DI
remains unchanged for this example. This observation alone does not establish invariance after general sample-adaptive reclustering or rematching.
